% =====================================================================
% PAPER I ("The Hammer")
% Exceptional eigenvalues of Bianchi orbifolds by verified finite
% elements. I: The method, and the spectral gap of the Picard orbifold
% ---------------------------------------------------------------------
% Sources: independent_exclusion/{PROOF,DESIGN,lower_bound_theory}.md
%          m3_certify.py, cr_prototype.py
% =====================================================================
\documentclass[11pt]{amsart}
\usepackage{amsmath,amssymb,amsthm,booktabs}
\newtheorem{theorem}{Theorem}[section]
\newtheorem{lemma}[theorem]{Lemma}
\newtheorem{proposition}[theorem]{Proposition}
\newtheorem{remark}[theorem]{Remark}
\newtheorem{corollary}[theorem]{Corollary}
\newcommand{\HH}{\mathbb{H}^3}
\newcommand{\ZZ}{\mathbb{Z}}
\newcommand{\RR}{\mathbb{R}}
\newcommand{\CC}{\mathbb{C}}
\newcommand{\PSL}{\mathrm{PSL}}
\newcommand{\vol}{\mathrm{vol}}

\title[Verified exclusion of exceptional eigenvalues, I]{Exceptional
eigenvalues of Bianchi orbifolds by verified finite elements.\\
I: The method, and the spectral gap of the Picard orbifold}
\author{}
\date{Draft, 2026-07-11}

\begin{document}
\begin{abstract}
We develop a machine-verified method for excluding Laplace eigenvalues
in $(0,1)$ on cusped hyperbolic $3$-orbifolds, combining a
Lax--Phillips cusp reduction, guaranteed Crouzeix--Raviart
finite-element lower bounds with fully explicit constants, and a
floating-point positive-definiteness certificate in the sense of Rump.
The method makes no use of the Selberg trace formula.  As a first
application we prove: the Picard orbifold
$\PSL(2,\ZZ[i])\backslash\HH$ has no Laplace eigenvalue in
$(0,1)$; equivalently $\lambda_1\ge 1$.  Every analytic constant is
enclosed in ball arithmetic and the final certificate is a
deterministic, reproducible computation.  Where certified Maass
computations establish existence of eigenvalues, the present work
establishes exclusion; the two approaches are complementary.  Papers~II
and~III apply the same machinery at congruence levels of the Picard
group and to the Eisenstein--Picard orbifold, where the statements are
new.
\end{abstract}
\maketitle

% ---------------------------------------------------------------------
\section{Introduction}

The Selberg eigenvalue problem asks for lower bounds on the first
positive Laplace eigenvalue $\lambda_1$ of arithmetic hyperbolic
quotients.  For congruence subgroups of $\mathrm{SL}(2,\ZZ)$ the
celebrated conjecture $\lambda_1\ge 1/4$ remains open in full
generality; the best unconditional bound is the Kim--Sarnak refinement
of the Selberg--Ramanujan conjecture.  Over imaginary quadratic fields
the analogous problem concerns Bianchi groups
$\Gamma\subset\PSL(2,\mathcal{O}_K)$ acting on hyperbolic $3$-space
$\HH$, with the Selberg-type threshold $\lambda_1\ge 1$.  The best
general unconditional bound is $\lambda_1\ge 975/1024$ of
Blomer--Brumley \cite{BB11}, obtained from the Kim--Sarnak $7/64$
exponent via the Jacquet--Langlands correspondence.  At level one for
the class-number-one fields $K=\mathbb{Q}(i)$ and $K=\mathbb{Q}(\omega)$,
positivity arguments based on the Selberg trace formula can push past
the threshold $B<1$ and exclude the entire exceptional interval $(0,1)$;
such an argument is implemented and certified for both groups in the
companion repository that houses the present work.  Those proofs rely
on the full conjugacy-class inventory, the scattering matrix, and a
carefully chosen test function.  Their margin degrades linearly with
volume, so they do not extend to congruence covers of moderate index.

The contribution of this paper is a \emph{different} architecture for
the same exclusion problem: a Lax--Phillips analysis of the cusp,
followed by a guaranteed finite-element lower bound on the compact core
and a verified positive-definiteness certificate.  The method never
invokes the Selberg trace formula.  It shares with the trace-formula
treatments only the classical fundamental domain.  As a validation
target we prove the level-one Picard statement, believed classically
known (cf.\ Elstrodt--Grunewald--Mennicke \cite{EGM98}, and Stramm's
small-eigenvalue bounds for Bianchi groups), by an independent route.

\begin{theorem}\label{thm:main}
The Laplace--Beltrami operator on
$L^2(\PSL(2,\ZZ[i])\backslash \HH)$ has no eigenvalue in
$(0,1)$.  Combined with the spectral decomposition of cofinite
Kleinian groups \cite[Thm.~3.8.1]{Friedman}, the spectrum in $(0,1)$ is
empty: $\lambda_1\ge 1$.
\end{theorem}

The architecture has three layers, each of independent interest.

\begin{enumerate}
\item \textbf{Reduction} (Section~\ref{sec:criterion}).
  A Lax--Phillips cusp analysis reduces the theorem to positivity of an
  explicit quadratic form $\mathcal{A}_s$ on a hyperplane in $H^1(K)$,
  where $K$ is the fundamental domain truncated at height $Y=1.25$.
  Nonzero Fourier modes in the cusp contribute nonnegatively; only the
  zero mode needs exact treatment.  Face identifications of the orbifold
  are relaxed away: testing over all of $H^1(K)$ is a valid relaxation.
\item \textbf{Guaranteed finite-element lower bound}
  (Sections~\ref{sec:mesh}--\ref{sec:G1}).
  A Crouzeix--Raviart discretization on an inner polyhedral mesh
  $K_h\subset K$ certifies the continuum pencil inequality from three
  \emph{finite} checks per spectral window: two scalar inequalities and
  one matrix positive-semidefiniteness check.  All interpolation and
  absorption constants are explicit.
\item \textbf{Verified certificate} (Section~\ref{sec:cert}).
  The finite checks are executed in ball arithmetic (Arb, 128-bit) for
  geometry and enclosures, with the matrix check certified by a literal
  implementation of Rump's positive-definiteness verification
  \cite{Rump06}.
\end{enumerate}

Related work falls into three strands.  First, rigorous Maass-form
computations (Booker--Str\"ombergsson--Venkatesh \cite{BSV06}, Child,
Seymour-Howell) establish \emph{existence} of eigenvalues to high
precision; exclusion requires a complementary architecture, which is
what we supply.  Second, the guaranteed FEM eigenvalue literature
(Carstensen--Gedicke \cite{CG14}, Liu, Carstensen--Puttkammer
\cite{CP22}, Carstensen--Zhai--Zhang \cite{CZZ20}, and the lower-energy
viewpoint of Carstensen--Dond--Maity--Nataraj \cite{CDMN26}) supplies the
template of projecting through a nonconforming interpolant and absorbing
errors by Young inequalities; our Theorem~G1 specializes this template
to the weighted hyperbolic forms and the pencil structure of the
exclusion criterion.  Third, verified numerical linear algebra in the
sense of Rump \cite{Rump06} converts a floating-point Cholesky run into
a rigorous positive-definiteness proof under IEEE-754 semantics.

Paper~II \cite{PaperII} applies the same engine to the Hecke congruence
quotients $\Gamma_0(\mathfrak{p})$ for Gaussian primes of norms $5$,
$9$, and $13$---statements unreachable by trace-formula positivity.
Paper~III \cite{PaperIII} treats the Eisenstein--Picard group
$\PSL(2,\ZZ[\omega])$ at level one and its first congruence rung
$\Gamma_0(1-\omega)$.  The present paper is the method paper: reviewer
burden is deliberately confined to (i)~the continuous-to-discrete
reduction, (ii)~the guaranteed FEM bounds, and (iii)~the interval
arithmetic.  All group-theoretic and combinatorial content is deferred.

A complete list of external mathematical inputs appears in
Section~\ref{sec:deps}.  Reproducibility commands are collected there as
well.  The certified run of Theorem~\ref{thm:main} is dated 2026-07-10.

% ---------------------------------------------------------------------
\section{The exclusion criterion}\label{sec:criterion}

\subsection{Setup}

Work in the upper half-space model
\[
  \HH = \bigl\{ P=(z,y):\; z=x_1+ix_2\in\CC,\; y>0\bigr\},
\]
with Riemannian metric $(|dz|^2+dy^2)/y^2$ and volume form
$dV=y^{-3}\,dx_1\,dx_2\,dy$.  For real $u$ the weighted Dirichlet and mass
forms are
\begin{align*}
  Q(u) &= \int |\nabla u|^2_{\mathrm{hyp}}\,dV
        = \int |\nabla_{\mathrm{euc}} u|^2\, y^{-1}\,dx\,dy,\\
  M(u) &= \int u^2\,dV.
\end{align*}
Let $\Gamma=\PSL(2,\ZZ[i])$.  The Humbert fundamental domain is
\[
  F=\bigl\{(z,y):\; x_1\in[-\tfrac12,\tfrac12],\;
  x_2\in[0,\tfrac12],\; |z|^2+y^2\ge 1\bigr\}.
\]
The cusp stabilizer is $\Gamma_\infty\cong\ZZ[i]\rtimes\langle z\mapsto -z\rangle$;
the cusp cross-section orbifold $T$ is the coordinate rectangle
$x_1\in[-\tfrac12,\tfrac12]$, $x_2\in[0,\tfrac12]$, of Euclidean area
$|T|=\tfrac12$.

Fix a truncation height $Y>1$ (we take $Y=1.25$; the strict inequality
avoids the mesh degeneracy where the unit sphere touches $y=1$ at $z=0$).
Split
\begin{align*}
  K &= F\cap\{y\le Y\}
    &&\text{(compact core: $y_f(z)\le y\le Y$, $y_f=\sqrt{1-|z|^2}$),}\\
  C &= F\cap\{y>Y\}
    &&\text{(cusp: $T\times(Y,\infty)$).}
\end{align*}

\begin{lemma}[Precisely invariant horoball; Shimizu]\label{lem:shimizu}
For $\gamma=\bigl(\begin{smallmatrix}a&b\\c&d\end{smallmatrix}\bigr)\in\Gamma$
with $c\neq 0$,
\[
  y(\gamma P)=\frac{y}{|cz+d|^2+|c|^2 y^2}\le\frac{1}{|c|^2 y}\le\frac{1}{y},
\]
since $|c|^2\ge 1$ for $0\neq c\in\ZZ[i]$.  Hence $\{y>Y\}$ with $Y\ge 1$
meets its $\Gamma$-translates only through $\Gamma_\infty$, and the cusp
Fourier analysis below is legitimate on the quotient.
\end{lemma}

The one-line proof above is self-contained (cf.\ Shimizu \cite{Shimizu63},
Leutbecher \cite{Leutbecher67}).

\subsection{Spectral preliminaries}

The spectral decomposition of $L^2(\Gamma\backslash\HH)$ for cofinite
Kleinian groups is recorded in Friedman \cite[Thm.~3.8.1]{Friedman}
(underlying eigenpacket theory: Elstrodt--Grunewald--Mennicke
\cite[§6.2]{EGM98}).  Every $f$ in the domain expands into discrete
eigenfunctions plus Eisenstein integrals $E_\alpha(\cdot,it)$.  Since
$E_\alpha(\cdot,it)$ has Laplace eigenvalue $1+t^2\ge 1$, the spectral
measure on $[0,1)$ is purely atomic, supported on the discrete
eigenvalues: \textbf{any spectral point in $(0,1)$ is an $L^2$
eigenvalue.}  Eigenfunctions are real-analytic (elliptic operator with
real-analytic coefficients; Morrey--Nirenberg), which also supplies unique
continuation.  The eigenvalue $\lambda=0$ has eigenspace equal to the
constants: $\Delta u=0$ with $u\in L^2$ on a connected finite-volume
quotient forces $\int|\nabla u|^2=0$.  Eigenfunctions with $\lambda\neq 0$
are orthogonal to constants by self-adjointness.

\subsection{Cusp Fourier analysis}

Let $u$ be a real eigenfunction, $\Delta u=\lambda u$, $\lambda=1-s^2$,
$s\in(0,1)$, $\|u\|=1$.  Work on the double cover of the cusp,
$\widetilde C=(\RR^2/\ZZ^2)\times(Y,\infty)$ (orbifold integrals equal
half the cover integrals).  Write Fourier modes $\hat u_\mu(y)$,
$\mu=(m,n)\in\ZZ^2$.

\begin{lemma}[Nonzero modes are harmless]\label{lem:nonzero}
For $\mu\neq 0$, the mode's contribution to $Q_C-\lambda M_C$ is
\begin{align*}
  &\tfrac12\int_Y^\infty\bigl(\hat u_\mu'^2+4\pi^2|\mu|^2\hat u_\mu^2\bigr)
    y^{-1}\,dy
  -\tfrac12\lambda\int_Y^\infty\hat u_\mu^2 y^{-3}\,dy\\
  &\qquad\ge\tfrac12\bigl(4\pi^2 Y^2-\lambda\bigr)
    \int_Y^\infty\hat u_\mu^2 y^{-3}\,dy\ge 0,
\end{align*}
using $\int\hat u^2 y^{-1}\ge Y^2\int\hat u^2 y^{-3}$ and
$4\pi^2 Y^2>1>\lambda$.  No boundary condition at $y=Y$ is needed.
\end{lemma}

\begin{lemma}[Exact zero-mode profile and defect]\label{lem:zero}
The zero mode $\hat u_0$ satisfies $y^2 f''-yf'+\lambda f=0$ on $(Y,\infty)$;
the solutions are $y^{1\pm s}$, and $L^2(y^{-3}\,dy)$ forces
$\hat u_0(y)=c\,(y/Y)^{1-s}$ with
$c=\hat u_0(Y)=(1/|T|)\int_T u(\cdot,Y)\,dx$.  Integration by parts with
the ODE gives \emph{exactly}
\[
  \int_Y^\infty\bigl(\hat u_0'^2-\lambda\hat u_0^2 y^{-2}\bigr)y^{-1}\,dy
  =-\hat u_0(Y)\hat u_0'(Y)/Y=-(1-s)c^2/Y^2,
\]
so the zero mode contributes $-(1-s)c^2/(2Y^2)$ to $Q_C-\lambda M_C$
(the factor $1/2$ from the cover).
\end{lemma}

\begin{lemma}[Orthogonality constraint]\label{lem:ortho}
From $\langle u,1\rangle=0$ and
$\int_Y^\infty(y/Y)^{1-s}y^{-3}\,dy=Y^{-2}/(1+s)$ one obtains
\[
  \int_K u\,dV+\frac{c}{2(1+s)Y^2}=0.
\]
\end{lemma}

\subsection{The criterion}

For $v\in H^1(K)$ define the Euclidean top-face integral
$t(v)=\int_T v(\cdot,Y)\,dx_1\,dx_2$ (so $c=2t$ on the cover), and for
$s\in(0,1)$, $\lambda=1-s^2$:
\begin{align*}
  \mathcal{A}_s(v)
    &= Q_K(v)-\lambda M_K(v)-2(1-s)\,t(v)^2/Y^2,\\
  \mathcal{L}_s(v)
    &= \int_K v\,dV+t(v)/\bigl((1+s)Y^2\bigr).
\end{align*}

\begin{theorem}[Exclusion criterion]\label{thm:criterion}
If for every $s\in(0,1)$,
\[
  \mathcal{A}_s(v)>0
  \quad\text{for all $0\neq v\in H^1(K)$ with $\mathcal{L}_s(v)=0$,}
\]
then $\Gamma\backslash\HH$ has no eigenvalue in $(0,1)$; i.e.\
$\lambda_1\ge 1$.
\end{theorem}

\begin{proof}
Let $u$ be an eigenfunction with $\lambda=1-s^2\in(0,1)$ and $\|u\|=1$.
Then
\[
  0=Q(u)-\lambda M(u)
  =[Q_K-\lambda M_K](u)+(\text{zero mode})+(\mu\neq 0\text{ modes})
  \ge\mathcal{A}_s(u|_K)
\]
by Lemmas~\ref{lem:nonzero}--\ref{lem:zero}, so $\mathcal{A}_s(u|_K)\le 0$.
Lemma~\ref{lem:ortho} gives $\mathcal{L}_s(u|_K)=0$.  Real-analyticity
(unique continuation) gives $u|_K\neq 0$.  Contradiction.
\end{proof}

\begin{remark}[Relaxation philosophy]
Testing over \emph{all} of $H^1(K)$ is a relaxation: no side or floor
identifications, no orbifold structure, plain Neumann-type free
boundaries.  If positivity survives this enlargement of the test space,
the certified computation never has to implement the face gluings.  The
criterion sees every $L^2$ eigenfunction---cuspidal or residual---exactly
as a $B<1$ positivity argument does.  Constants sanity check:
$v\equiv 1$ has $\mathcal{A}_s<0$ but $\mathcal{L}_s(1)\neq 0$, so the
constraint is what excludes the $\lambda=0$ eigenfunction, as it must.
For fixed $v$, $\mathcal{A}_s(v)$ is decreasing in $\lambda$, so a finite
$\lambda$-grid checked at right endpoints controls each window; the
$s$-dependence of the constraint is handled by the interval-$\kappa$
reduction of Section~\ref{sec:G1}.
\end{remark}

\subsection{Pencil form}

The constrained positivity of Theorem~\ref{thm:criterion} is equivalent
to an unconstrained pencil inequality, which is the natural container for
nonconforming lower bounds.

\begin{lemma}[Pencil criterion]\label{lem:P}
Fix $\rho>0$ and $\varepsilon_0>0$.  If for every $s\in(0,1)$
\begin{equation}\label{eq:Ps}
  Q(v)+\rho\,\mathcal{L}_s(v)^2
  \ge(1+\varepsilon_0)\bigl[\lambda M(v)+\beta(s)\,t(v)^2\bigr]
  \qquad\text{for all $v\in H^1(K)$,}
\end{equation}
where $\lambda=1-s^2$ and $\beta(s)=2(1-s)/Y^2$, then the criterion of
Theorem~\ref{thm:criterion} holds (with strict positivity), hence
$\lambda_1\ge 1$.
\end{lemma}

\begin{proof}
Let $v\neq 0$ with $\mathcal{L}_s(v)=0$.  Then \eqref{eq:Ps} gives
$\mathcal{A}_s(v)=Q-\lambda M-\beta t^2
\ge\varepsilon_0(\lambda M(v)+\beta t(v)^2)\ge\varepsilon_0\lambda M(v)>0$.
\end{proof}

The free parameter $\rho$ never affects correctness; it only affects the
numerical margin of the discrete checks.

% ---------------------------------------------------------------------
\section{Geometry: the inner polyhedral mesh}\label{sec:mesh}

Start from a mapped hexahedral grid of $N_1\times N_2\times N_3$ cells over
$(\xi_1,\xi_2,\xi_3)\in[-\tfrac12,\tfrac12]\times[0,\tfrac12]\times[0,1]$,
with vertical coordinate $y=y_f(1-\xi_3)+Y\xi_3$.  \textbf{Floor lift:}
every node with $\xi_3=0$ is placed at $y=y_f(z)+\ell$ instead of $y_f(z)$,
with the lift
\[
  \ell:=\tfrac12\cdot\max_{\mathrm{cell}}|D^2 y_f|\cdot
  (\text{cell diameter in $z$})^2,\qquad
  |D^2 y_f|\le y_f^{-3}\le 2^{3/2},
\]
evaluated per floor cell.  Each (possibly non-planar) hexahedral cell is
split into six tetrahedra by the Kuhn triangulation with globally
consistent diagonals; all tet faces are planar and the mesh is
face-to-face.  Write $K_h$ for the union of the tets, $\Gamma_h$ for its
floor (a piecewise-linear graph $\hat y(z)$ over the rectangle
$R=[-\tfrac12,\tfrac12]\times[0,\tfrac12]$), and
$\Sigma:=K\setminus K_h$ for the sliver between the polyhedral floor and
the sphere.  Set $\delta(z):=\hat y-y_f\ge 0$.

\begin{lemma}[Inclusion]\label{lem:G}
The function $y_f$ is concave on $R$.  On a floor triangle with longest
$z$-edge $d$, the chord interpolant lies below $y_f$ by at most the sag
$(\max|D^2 y_f|)\cdot d^2/8$.  With the per-cell lift
$\ell:=(\max_{\mathrm{cell}} y_f^{-3})\cdot d^2/8$ assigned to nodes as
the maximum over adjacent cells, one has $\hat y\ge y_f$ pointwise, i.e.\
$K_h\subset K$, and $\delta(z)\le\max$ vertex lift.
\end{lemma}

\begin{proof}[Sketch]
Taylor with concavity: the chord at $z=\sum\theta_i z_i$ satisfies
$\mathrm{chord}(z)\ge y_f(z)-\tfrac12\max|D^2 y_f|\sum\theta_i|z_i-z|^2$.
The barycentric variance $\sum\theta_i|z_i-z|^2$ is at most $d^2/4$
(maximized at an edge midpoint of the longest edge).  Adding the lifts
raises the chord by at least the sag.  Since $\mathrm{chord}\le y_f$
without lift, also $\delta\le\max$ vertex lift.
\end{proof}

\begin{remark}[Mesh-split agnosticism]
The level-one certificate of Theorem~\ref{thm:main} uses a Kuhn split.
Paper~II uses the symmetric $24$-split (body centre plus face centres),
required for mirror-compatible gluing at congruence level.  Lemmas~G, E,
S, I1 and Theorem~G1 below do not depend on the split.
\end{remark}

Per-tet data (all computed by code): diameter $h_T$, volume $|T|$, face
areas; $y$-range $[y_T^-,y_T^+]$ over $T$ (min/max over the four vertices,
valid since $y$ is affine on $T$); sandwiched weights
\[
  w_T^Q:=(y_T^+)^{-1}\quad\text{(stiffness, from below),}\qquad
  w_T^M:=(y_T^-)^{-3}\quad\text{(mass, from above),}
\]
so that $Q(v)\ge Q_{\mathrm{pc}}(v):=\sum_T w_T^Q\|\nabla v\|_{L^2(T)}^2$ and
$M|_{K_h}(v)\le M_{\mathrm{pc}}(v):=\sum_T w_T^M\|v\|_{L^2(T)}^2$ for every
$v\in H^1(K_h)$.

% ---------------------------------------------------------------------
\section{The guaranteed lower bound}\label{sec:G1}

This section turns the infinite-dimensional positivity of
Lemma~\ref{lem:P} into a finite, verified computation.  Every constant is
explicit and interval-computable.

\subsection{Crouzeix--Raviart interpolation}

Let $V_h:=\mathrm{CR}^1(\mathcal{T})$ be the space of piecewise affine
functions continuous at face barycentres (equivalently: matching face
means), with \emph{no} boundary conditions.  The face-mean interpolant
$I=I_{\mathrm{CR}}:H^1(K_h)\to V_h$ satisfies, on each tet $T$, with
$e:=v-Iv$:
\begin{align*}
  \text{(I2)}\quad
  &\nabla(Iv)|_T=\Pi_0(\nabla v|_T)
  \;\Rightarrow\;
  \int_T\nabla e\cdot\nabla w_h=0\text{ for }w_h\in V_h,\\
  &\text{hence }
  \|\nabla v\|_T^2=\|\nabla Iv\|_T^2+\|\nabla e\|_T^2
  \text{ and }
  Q_{\mathrm{pc}}(v)=Q_{\mathrm{pc}}(Iv)+Q_{\mathrm{pc}}(e);\\
  \text{(I1)}\quad
  &\|e\|_{L^2(T)}\le\kappa\, h_T\|\nabla e\|_{L^2(T)}.
\end{align*}
Property (I2) is immediate from the face-mean property
($\int_T\nabla e=\sum_F\nu_F\int_F e=0$ and $\nabla Iv$ is constant).  For
(I1) we use a self-contained constant.

\begin{lemma}[CR interpolation on arbitrary tetrahedra]\label{lem:I1}
Inequality (I1) holds with
\[
  \kappa_{\mathrm{sc}}:=\sqrt{\frac1{\pi^2}+\frac1{15}}\le 0.40988.
\]
\end{lemma}

\begin{proof}[Sketch]
The error $e=v-I_{\mathrm{CR}}v$ has vanishing face means.  Split
$e=(e-\bar e_T)+\bar e_T$, orthogonal in $L^2(T)$.  Payne--Weinberger for
the convex domain $T$ \cite{PW60,Bebendorf03} gives
$\|e-\bar e_T\|\le(h_T/\pi)\|\nabla e\|$.  For the mean, pick a vertex $P$
with opposite face $F$: the divergence identity
$\operatorname{div}((x-P)e)=3e+(x-P)\cdot\nabla e$ and the face-mean
vanishing yield $3|T|\bar e_T=-\int_T(x-P)\cdot\nabla e$; Cauchy--Schwarz
and the barycentric second-moment bound $m_2\le(3/5)h_T^2$ give
$|T|\bar e_T^2\le(h_T^2/15)\|\nabla e\|^2$.  Combining,
$\kappa^2\le 1/\pi^2+1/15$.
\end{proof}

\begin{remark}
The sharper published constant
$\kappa_1:=\sqrt{1/\pi^2+1/120}\le 0.33115$ of Carstensen--Puttkammer
\cite[§2.4]{CP22} (sourced to Carstensen--Zhai--Zhang \cite{CZZ20}) is
available as an optional sharpening.  The certificate runs with
$\kappa_{\mathrm{sc}}$ by default; nothing below depends on which is
chosen, since $\kappa$ enters only scalar constants.
\end{remark}

\begin{lemma}[Exact trace reproduction]\label{lem:D0}
The top plane $\{y=Y\}$ is exactly tiled by boundary faces of the mesh,
and $I_{\mathrm{CR}}$ matches the mean of $v$ on every mesh face.  Hence
$t(e)=0$ identically: the trace functional is exactly reproduced, no
slab bound is needed, and the Young split of $t(v)^2$ in the proof of
Theorem~\ref{thm:G1} can be skipped (boundary term uninflated in $D_h$).
\end{lemma}

\begin{remark}
Historically the level-one certificate used a positive slab constant
$\tau_h>0$, which is valid but lossier; the certified run therefore
stands a fortiori under Lemma~\ref{lem:D0}.  The exact-reproduction
viewpoint became load-bearing in Paper~II, where multi-cusp boundary
blocks must remain uninflated.
\end{remark}

\subsection{Error and sliver lemmas}

Write $a(v)=\int_K v\,y^{-3}\,dx\,dy$ and
$\kappa_c(s)=1/((1+s)Y^2)$, so
$\mathcal{L}_s(v)=a(v)+\kappa_c(s)\,t(v)$.

\begin{lemma}[Functional errors]\label{lem:E}
For any $v\in H^1(K_h)$ and $e=v-I_{\mathrm{CR}}v$:
\begin{align*}
  \text{(E-M)}\quad
  &M_{\mathrm{pc}}(e)\le\gamma^2 Q_{\mathrm{pc}}(e),
  \quad\gamma:=\kappa\cdot\max_T h_T\sqrt{w_T^M/w_T^Q};\\
  \text{(E-a)}\quad
  &|a(e)|\le\alpha_h\sqrt{Q_{\mathrm{pc}}(e)},
  \quad\alpha_h:=\kappa\Bigl(\sum_T(w_T^M)^2|T|h_T^2/w_T^Q\Bigr)^{1/2};\\
  \text{(E-$\mathcal{L}$)}\quad
  &|\mathcal{L}_{s_0}(e)|\le\sigma_h\sqrt{Q_{\mathrm{pc}}(e)},
  \quad\sigma_h:=\alpha_h+\kappa_c(s_0)\,\tau_h,
\end{align*}
where $\tau_h=0$ under Lemma~\ref{lem:D0}, and otherwise $\tau_h$ is the
slab-mean constant of the historical route (any slab height
$0<H_t\le Y-\max\hat y$).
\end{lemma}

\begin{lemma}[Sliver via first-layer column means]\label{lem:S}
Fix a column height $H_s\in(0,Y-\max\hat y]$.  For each floor face $F$
let $\pi(F)\subset R$ be its $z$-projection and $\hat\delta_F$ the max
vertex lift among its vertices.  The columns
$\mathrm{col}(F):=\{(z,y):z\in\pi(F),\;\hat y(z)\le y\le\hat y(z)+H_s\}$
are disjoint subsets of $K_h$.  Then for every $v\in H^1(K)$,
\begin{align*}
  M_\Sigma(v)&:=\int_\Sigma v^2 y^{-3}
  \le S_M M_{K_h}(v)+S_Q Q_{K_h}(v)+S_\Sigma Q_\Sigma(v),\\
  |a_\Sigma(v)|&:=\bigl|\int_\Sigma v\,y^{-3}\bigr|
  \le\sqrt{V_\Sigma}\sqrt{M_\Sigma(v)},
\end{align*}
with explicit constants $S_M,S_Q,S_\Sigma,V_\Sigma$ computed from per-face
heights and $y$-ranges (formulas in the companion theory note; all are
$O(h)$ or $O(h^2)$).
\end{lemma}

\subsection{Exact main terms}

The Young split $v=Iv+e$ is pointwise, so the \emph{main} terms need no
weight freezing---only the error steps (E-M) and (E-a), which convert
$\|e\|$-type norms to $\|\nabla e\|$, use sandwiched weights.  Define per
tet the nonnegative stiffness remainder
\[
  Q_{\mathrm{rem}}:=\sum_T\Bigl(\int_T y^{-1}-w_T^Q|T|\Bigr)\nabla\varphi\nabla\varphi^T
  \succeq 0
\]
(with $\int_T y^{-1}$ Taylor-enclosed), the exact-weight mass $M_{\mathrm{ex}}$
(entries $\int_T\varphi_a\varphi_b y^{-3}$, Taylor-enclosed like $\ell_0$),
and $\rho_w:=\max_T(y_T^+/y_T^--1)$.  Then for any $\theta_4\in(0,1)$:
\begin{align*}
  Q(v)&\ge\bigl[Q_{\mathrm{pc}}+(1-\theta_4)Q_{\mathrm{rem}}\bigr](Iv)
    +\bigl(1-(1/\theta_4-1)\rho_w\bigr)Q_{\mathrm{pc}}(e),\\
  M(v)&\le(1+\alpha)M_{\mathrm{ex}}(Iv)+(1+1/\alpha)\gamma^2 Q_{\mathrm{pc}}(e).
\end{align*}
(The pure pc-sandwich alone costs roughly $38\%$ mass inflation on the
level-one mesh; the exact-main-term refinement recovers this.)

\subsection{Master theorem}

Fix a window $W=[s^-,s^+]\subseteq[0,1]$ and a reference $s_0\in W$.
Worst-case data:
\[
  \lambda^+:=1-(s^-)^2,\qquad
  \beta^+:=2(1-s^-)/Y^2,\qquad
  \Delta\kappa:=\max\bigl(|\kappa_c(s^-)-\kappa_c(s_0)|,\,
  |\kappa_c(s^+)-\kappa_c(s_0)|\bigr).
\]
Choose Young parameters $\theta,\theta_2,\theta',\alpha\in(0,1)$, a target
$\nu^*=1+\varepsilon_0$, and $\rho>0$.  Define, in this order (all
computable from mesh data):
\begin{align*}
  \omega&:=2\rho(1/\theta_2-1)V_\Sigma,\\
  c_Q&:=1-(\omega+\nu^*\lambda^+)S_Q,\\
  c_\Sigma&:=1-(\omega+\nu^*\lambda^+)S_\Sigma
    \quad\text{(must be $\ge 0$; then $Q_\Sigma$ is dropped),}\\
  \tilde\lambda&:=\nu^*\lambda^+(1+S_M)+\omega S_M,\\
  \tilde\beta&:=\nu^*\beta^++2\rho(1/\theta_2-1)\Delta\kappa^2,\\
  \tilde\rho&:=\rho(1-\theta_2),\\
  c_e&:=c_Q\bigl(1-(1/\theta_4-1)\rho_w\bigr)
    -\tilde\rho(1/\theta-1)\sigma_h^2,\\
  d_e&:=\tilde\lambda(1+1/\alpha)\gamma^2
    +\tilde\beta(1+1/\theta')\tau_h^2.
\end{align*}
(Under Lemma~\ref{lem:D0} one has $\tau_h=0$ and the second summand of
$d_e$ vanishes; the historical level-one run retained $\tau_h>0$.)
Discrete check: with the CR matrices $Q_{\mathrm{pc}}$, $M_{\mathrm{pc}}$
(or $M_{\mathrm{ex}}$), vectors $\ell_0$ (for
$\mathcal{L}_{s_0}(v_h)=a(v_h)+\kappa_c(s_0)t(v_h)$) and $t$,
\begin{align*}
  N_h&:=c_Q\bigl[Q_{\mathrm{pc}}+(1-\theta_4)Q_{\mathrm{rem}}\bigr]
    +\tilde\rho(1-\theta)\,\ell_0\ell_0^T,\\
  D_h&:=\tilde\lambda(1+\alpha)M_{\mathrm{ex}}
    +\tilde\beta(1+\theta')\,tt^T.
\end{align*}

\begin{theorem}[Master lower bound]\label{thm:G1}
Suppose $c_\Sigma\ge 0$, $c_e>0$, and
\begin{enumerate}
\item[(a)] $N_h-D_h\succeq 0$ (positive semidefinite on the CR space),
\item[(b)] $c_e/d_e\ge 1$.
\end{enumerate}
Then the pencil inequality \eqref{eq:Ps} holds with $1+\varepsilon_0=\nu^*$
for every $s\in W$.  If (a) and (b) hold for a family of windows covering
$(0,1)$, then by Lemma~\ref{lem:P} one has $\lambda_1\ge 1$.
\end{theorem}

\begin{proof}[Outline]
Let $v\in H^1(K)$ and $s\in W$.  Write $N:=Q(v)+\rho\mathcal{L}_s(v)^2$ and
$D:=\lambda M(v)+\beta(s)t(v)^2$.  Since $\lambda\le\lambda^+$ and
$\beta(s)\le\beta^+$, one has $D\le\lambda^+ M(v)+\beta^+ t(v)^2$.

\emph{Step 1} (window shift + sliver in $\mathcal{L}$).
$\mathcal{L}_s(v)=\mathcal{L}_{s_0}(v)|_{K_h}+a_\Sigma(v)
+(\kappa_c(s)-\kappa_c(s_0))t(v)$; Young and Lemma~\ref{lem:S} absorb the
error into $\omega$ and $\tilde\beta$.

\emph{Step 2} (sliver in $M$).  $M(v)=M_{K_h}(v)+M_\Sigma(v)$ and
Lemma~\ref{lem:S} produce the effective mass coefficient $\tilde\lambda$
and the nonnegativity condition $c_\Sigma\ge 0$ that drops $Q_\Sigma$.

\emph{Step 3} (pc sandwich and exact main terms).
$Q_{K_h}\ge Q_{\mathrm{pc}}+(1-\theta_4)Q_{\mathrm{rem}}$ on the main term
and $M_{K_h}\le M_{\mathrm{ex}}$ (or $M_{\mathrm{pc}}$) reduce the continuum
comparison to discrete data.

\emph{Step 4} (CR projection).  Split $v=Iv+e$.  By (I2),
$Q_{\mathrm{pc}}(v)=Q_{\mathrm{pc}}(Iv)+Q_{\mathrm{pc}}(e)$.  Lemma~\ref{lem:E}
and Young give
\begin{align*}
  \mathcal{L}_0(v)^2&\ge(1-\theta)\mathcal{L}_0(Iv)^2
    -(1/\theta-1)\sigma_h^2 Q_{\mathrm{pc}}(e),\\
  M(v)&\le(1+\alpha)M_{\mathrm{ex}}(Iv)
    +(1+1/\alpha)\gamma^2 Q_{\mathrm{pc}}(e),\\
  t(v)^2&\le(1+\theta')t(Iv)^2+(1+1/\theta')\tau_h^2 Q_{\mathrm{pc}}(e).
\end{align*}
The bracketed main terms satisfy the discrete comparison by (a); the
error terms satisfy $c_e Q_{\mathrm{pc}}(e)\ge d_e Q_{\mathrm{pc}}(e)$ by
(b).  Adding proves the continuum inequality on $K_h$, hence on $K$.
\end{proof}

\begin{remark}[Efficiency]
All corrections are $O(h)$ or $O(h^2)$: $\gamma,\alpha_h=O(h)$,
$\delta,S_*,V_\Sigma=O(h^2)$.  Hence $c_e/d_e=O(h^{-2})\to\infty$ and
$N_h-D_h$ tends to the unperturbed discrete pencil.  The theorem is not
just true but efficient: the $h$-dependence enters only through small
multiplicative and additive corrections.  What the code must certify per
window is $c_\Sigma\ge 0$, $c_e>d_e$, and $N_h-D_h\succeq 0$---no eigenvalue
\emph{solve} is part of the certificate, only positive-definiteness
checks.
\end{remark}

% ---------------------------------------------------------------------
\section{The verified certificate}\label{sec:cert}

\subsection{Exactness conventions}

Mesh coordinates are IEEE doubles representing exact dyadic rationals
from the mapped grid.  Geometry admissibility (Lemma~\ref{lem:G}:
$\hat y\ge y_f$ on every floor face) is re-verified in ball arithmetic.
Per-tet data and sandwiched weights are interval-enclosed.  The only
curved-weight objects requiring nontrivial enclosures are the vector
$\ell_0$ (entries $\int_T\varphi_F y^{-3}$) and the exact-weight mass and
stiffness-remainder moments.

\subsection{Interval enclosures}

Integrals of the form $\int_T\varphi_a y^{-3}\,dx$ are enclosed by a
degree-$p$ Taylor expansion of $y^{-3}$ about the tet mean $\bar y$: the
moments $\int_T\varphi_a(y-\bar y)^k$ are exact barycentric-monomial
integrals
$\bigl(\int_T\prod\lambda_i^{\alpha_i}=6|T|\prod\alpha_i!/(|\alpha|+3)!\bigr)$,
and the remainder is a ball controlled by
$y_T^{-(4+p)}\Delta^{p+1}$.  With $p=5$ the vector radius is of order
$10^{-7}$ against eigenvalue margins of order $10^{-4}$.  Consistency
checks include row sums of $M_{\mathrm{ex}}$ against $\ell$ and the
identity $1^TM1=\vol(K)$ up to the $O(h^2)$ sliver.

\begin{lemma}[Factored interval rank-one reduction]\label{lem:R}
Let $\ell\in\RR^n$ lie componentwise in $[\hat\ell\pm r]$.  Then
\[
  \ell\ell^T\succeq\hat\ell\hat\ell^T-2\|\hat\ell\|_2\|r\|_2\cdot I
  \quad\text{and}\quad
  \ell\ell^T\preceq\hat\ell\hat\ell^T+(2\|\hat\ell\|\|r\|+\|r\|^2)\cdot I.
\]
\end{lemma}

\begin{proof}
Write $\ell=\hat\ell+\delta$ with $|\delta^Tx|\le\|r\|\|x\|$ and
$|\hat\ell^Tx|\le\|\hat\ell\|\|x\|$.
\end{proof}

\subsection{Rump's certificate}

The matrix check $N_h-D_h\succeq 0$ for all interval data reduces, by
Lemma~\ref{lem:R} and entrywise radius domination, to
$\lambda_{\min}(\hat A)>\varepsilon_{\mathrm{tot}}$ for a float midpoint
matrix $\hat A$ and an explicit total shift $\varepsilon_{\mathrm{tot}}$.
This is certified by a literal implementation of Rump
\cite{Rump06}, Theorems~2.3, Corollary~2.4, and Lemma~2.5:

\begin{itemize}
\item \emph{Theorem 2.3.}  If the floating-point Cholesky factorization of
  a symmetric $A\in\mathbb{F}^{n\times n}$ runs to completion, then
  $\lambda_{\min}(A)>-\|\Delta(A)\|_2$, where $\Delta(A)_{ij}$ is built from
  $\gamma_{s(i,j)+2}$, the diagonal scaling $d_j$, and the underflow term
  $M\cdot\eta$ (IEEE-754 double: $\varepsilon=2^{-53}$, $\eta=2^{-1074}$).
\item \emph{Corollary 2.4 + Lemma 2.5.}  With $c\ge\|\Delta(A)\|_2
  +\varepsilon_{\mathrm{tot}}$ in $\mathbb{F}$, form the $\varphi$-trick
  diagonal $\tilde a_{ii}=\mathrm{fl}(d_i'-\varphi|d_i'|)$ guaranteeing
  $\tilde a_{ii}\le a_{ii}-c$ without directed rounding.  If floating
  Cholesky of $\tilde A$ runs to completion, then
  $\lambda_{\min}(A)\ge\varepsilon_{\mathrm{tot}}$.
\item Library routine: Rump notes that ``any library routine can be used'';
  the code calls LAPACK \texttt{dpotrf} (as in INTLAB \texttt{isspd}).
\end{itemize}

\subsection{Two refinements (machinery of the Hammer)}

Two refinements, stress-tested into existence by the $N\mathfrak{p}=13$
certificate of Paper~II and belonging to the method rather than any single
theorem, are recorded here.

\begin{proposition}[Error-free diagonal congruence]\label{prop:SAS}
The scaling $A\mapsto SAS$ with
$s_i=2^{\mathrm{round}(-\frac12\log_2 a_{ii})}$ is exact in IEEE-754
(powers of two), preserves definiteness, and equilibrates the diagonal.
It removes the mismatch between Rump's stiffness-scale shift and the
mass-scale eigenvalue floor.
\end{proposition}

\begin{proposition}[Per-row radius absorption]\label{prop:perrow}
For symmetric $|E|\le R$ entrywise,
$E\succeq-\mathrm{diag}(r)$ with $r_i=\sum_j R_{ij}$.  This replaces the
uniform max-row-sum shift.  Proof: $2|x_i x_j|\le x_i^2+x_j^2$.
\end{proposition}

At the level-one size ($n\approx 5\cdot 10^3$) the unscaled Rump path
already succeeds; at $n\approx 3\cdot 10^4$ (Paper~II, $N\mathfrak{p}=13$)
the uniform shift exceeds $\lambda_{\min}(A)\sim 10^{-5}$ while the matrix
is still SPD, and the two refinements restore the certificate.

\subsection{Frozen certificate for Theorem~\ref{thm:main}}

\begin{table}[ht]
\centering
\begin{tabular}{@{}ll@{}}
\toprule
Parameter & Value \\
\midrule
Mesh & $12\times 6\times 6$ mapped hex cells $\to$ $2592$ tets, $5544$ CR dofs \\
$Y$ & $1.25$ \\
$(\rho,\theta,\theta_2,\theta',\alpha)$ & $(55,\,0.7,\,0.5,\,0.5,\,0.5)$ \\
$\nu^*$ & $1.05$ \\
$s$-grid & $8$ uniform windows covering $[0,1]$ \\
$\ell_0$ Taylor order & $p=5$ \\
Arb precision & $128$ bits \\
$\kappa$ & $\kappa_{\mathrm{sc}}=\sqrt{1/\pi^2+1/15}$ (also passes with CZZ) \\
\bottomrule
\end{tabular}
\caption{Frozen certificate parameters for Theorem~\ref{thm:main}.}
\label{tab:frozen1}
\end{table}

Certified run (2026-07-10), abridged: geometry admissible
(Lemma~\ref{lem:G}, arb-verified); constants (arb, upper bounds)
$\gamma\approx 0.09847$, $\tau\approx 0.29085$, $\alpha_h\approx 0.02240$,
$S_M\approx 1.34\cdot 10^{-1}$, $S_Q\approx 7.45\cdot 10^{-3}$,
$S_\Sigma\approx 5.71\cdot 10^{-3}$, $V_\Sigma\approx 6.94\cdot 10^{-3}$;
radii $\|\ell_{\mathrm{rad}}\|\approx 6.6\cdot 10^{-8}$,
$\|Q_r\|\approx 2.5\cdot 10^{-13}$; all eight windows:
$c_e>d_e$ true, PSD true (Rump shift $c\approx 2\cdot 10^{-8}$, float
diagnostic margins $\approx 10^{-4}$).

The parameter-search discipline is: floating-point experiments choose
mesh, Young parameters, and $\rho$; the frozen tuple is then verified in
arb.  Floats choose; arb verifies.  Correctness never depends on the
search path.

% ---------------------------------------------------------------------
\section{Validation and diagnostics}\label{sec:valid}

A conforming Q1 prototype on the mapped box
(trilinear elements, dense assembly of $S$, $M$, $a$, $t$) measures the
margin
\[
  \mu(\lambda):=\min\bigl\{\mathcal{A}_s(v):\;\mathcal{L}_s(v)=0,\;
  M(v)=1\bigr\}
\]
as a constrained generalized eigenvalue.  With $Y=1.25$ on meshes
$16\times 8\times 6$ and $24\times 12\times 9$ (agreeing to three digits):
\[
  \mu(\lambda)\text{ decreases from }\approx 8.23\ (\lambda=0.05)
  \text{ to }\approx 7.28\ (\lambda=0.999);\quad
  \min\mu\approx 7.3.
\]
The minimizer is the first nonconstant Neumann mode of the core
($\lambda_2^N\approx 8.28$), which has $t(v)\approx 0$ and mean $\approx 0$
by parity, so $\mu(\lambda)\approx\lambda_2^N-\lambda$.  Consequences:
the Neumann relaxation is comfortably enough; the certification target
is modest (any guaranteed lower bound $\ge\lambda+\varepsilon$ for the
constrained problem suffices, with roughly seven units of headroom for CR
constants, weight sandwiching, and interval slop); without the
constraint the near-constant direction gives $\approx-3.4$, confirming
that the boundary term and constraint are load-bearing and correctly
signed.

A Neumann box model (exact first eigenvalue known) validates the
guaranteed-lower-bound pipeline: the discrete GLB sits below the exact
value, which sits below the conforming Galerkin upper bound.  The CR
float pipeline passes all eight $s$-windows at $\nu^*=1.05$ on the
$12\times 6\times 6$ mesh, with $c_e/d_e\ge 1.88$.  Robustness holds under
both $\kappa_{\mathrm{sc}}$ and the sharper CZZ constant.

The parent repository certifies $B(\ZZ[i])\le 0.3199<1$ via the trace
formula, hence the same conclusion $\lambda_1\ge 1$.  The two proofs share
only the fundamental domain.  Agreement of two disjoint methods on the
level-one statement is the independent validation that motivated the
present architecture---and the architecture is reusable at congruence
levels where the statement is new (Paper~II).

% ---------------------------------------------------------------------
\section{Dependencies and reproducibility}\label{sec:deps}

\subsection{External mathematical inputs}

\begin{enumerate}
\item \textbf{Humbert fundamental domain} for $\PSL(2,\ZZ[i])$
  (classical; the only input shared with the companion
  trace-formula engine).
\item \textbf{Shimizu's lemma} (precisely invariant horoball)---one-line
  self-contained proof as Lemma~\ref{lem:shimizu}
  (cf.\ \cite{Shimizu63,Leutbecher67}).
\item \textbf{Spectral decomposition of cofinite Kleinian groups}---
  Friedman \cite[Thm.~3.8.1]{Friedman}, resting on
  \cite[§6.2]{EGM98}.  Used only for: (i)~spectrum $\cap[0,1)$ is atomic
  (Eisenstein part has eigenvalue $1+t^2\ge 1$); (ii)~the $\lambda_1\ge 1$
  phrasing.  The core statement ``no $L^2$ eigenvalues in $(0,1)$'' needs
  neither.
\item \textbf{Real-analyticity of eigenfunctions} (Morrey--Nirenberg;
  supplies unique continuation).
\item \textbf{Payne--Weinberger inequality} for convex domains
  \cite{PW60,Bebendorf03}---the only external input to Lemma~\ref{lem:I1}.
\item \textbf{Rump, verification of positive definiteness}
  \cite{Rump06}---Thm.~2.3, Cor.~2.4, Lemma~2.5.  Requires IEEE-754
  semantics of the floating Cholesky; ``any library routine can be used''.
\end{enumerate}

Everything else---cusp Fourier analysis, Hardy/ODE lemmas, the trace
identity, Lemmas~G, S, E, I1, R, and Theorem~G1---is proved from scratch
in the companion theory notes and sketched above.

\subsection{Explicitly not used}

The Selberg trace formula; any test function $h$; the scattering matrix
or Dedekind zeta; the torsion or conjugacy-class inventories; Then's
computed spectrum; and every result file of the companion trace-formula
engine.

\subsection{Software}

Correctness-critical: python-flint (Arb ball arithmetic).  Covered by
Rump's theorem: NumPy/SciPy/LAPACK \texttt{dpotrf}.  Everything else is
diagnostic.  IEEE-754 double arithmetic is assumed throughout the Rump
path.

\subsection{Reproduce}

From the directory \texttt{independent\_exclusion/}:
\begin{verbatim}
python cr_prototype.py     # float pipeline + box-model validation
python m3_certify.py       # the certificate (both kappa modes)
\end{verbatim}
Both $\kappa$ modes (Lemma~I1 default and optional CZZ) pass all eight
windows.

% ---------------------------------------------------------------------
\section{Outlook}

The measured float margin $\mu\approx 7.3$ at level one leaves comfortable
headroom for the discrete checks.  Paper~II climbs the congruence ladder
$\Gamma_0(\mathfrak{p})$ for Gaussian primes of norms $5$, $9$, and $13$,
with measured margins $\mu\approx 4.4$, $2.4$, and $1.45$ respectively---
sublinear degradation against volume, in contrast to the linear
degradation of the trace-formula constant $B$.  Those theorems are new:
trace-formula positivity already fails at $N\mathfrak{p}=5$ ($B\approx 1.9$).
Paper~III carries the method to $\PSL(2,\ZZ[\omega])$ at level one and to
$\Gamma_0(1-\omega)$.

Complementary existence tools (interval-Hejhal, rigorous Maass
computations) remain the right instruments for producing eigenvalues
above the gap; the present work produces rigorous absences below it.
The combination---certified exclusion on $(0,1)$ together with certified
existence of the first Maass form---gives a complete rigorous picture of
the bottom of the spectrum for the quotients within reach of both
methods.

% ---------------------------------------------------------------------
\appendix
\section{Deferred proof details}

\subsection*{Trace identity}

\begin{lemma}[Trace identity]\label{lem:T}
Let $T$ be a tetrahedron, $F$ a face, and $P$ the opposite vertex.  For
any $w\in H^1(T)$:
\[
  \int_F w^2\,d\sigma
  =\frac{|F|}{|T|}\int_T\Bigl[w^2+\tfrac23 w(x-P)\cdot\nabla w\Bigr]dx
  \le\frac{|F|}{|T|}\bigl(\|w\|_T^2+\tfrac23 h_T\|w\|_T\|\nabla w\|_T\bigr).
\]
\end{lemma}

\begin{proof}
$\operatorname{div}((x-P)w^2)=3w^2+2w(x-P)\cdot\nabla w$; integrate over
$T$.  On the three faces containing $P$, $(x-P)\cdot\nu=0$; on $F$,
$(x-P)\cdot\nu=\operatorname{dist}(P,\mathrm{aff}\,F)=3|T|/|F|$.
Cauchy--Schwarz and $|x-P|\le h_T$ finish the bound.
\end{proof}

This identity is used in the historical per-face route to $\tau_h$ and as
a tool inside Lemma~\ref{lem:I1}.

\subsection*{Historical slab form of (E-t)}

For the record: before Lemma~\ref{lem:D0} was recognized, the trace error
was controlled by a slab-mean identity.  For a.e.\ $z\in R$ and
$w\in H^1$,
\[
  w(z,Y)=\frac1{H_t}\int_{Y-H_t}^Y w\,dy
  +\frac1{H_t}\int_{Y-H_t}^Y(y-Y+H_t)\partial_y w\,dy,
\]
valid whenever the slab $\{Y-H_t\le y\le Y\}\subset K_h$.  Cauchy--Schwarz
and the CR estimate on the slab tets produce
\[
  |t(e)|\le\tau_h\sqrt{Q_{\mathrm{pc}}(e)},\qquad
  \tau_h=\sqrt{\frac{|R|}{H_t}}\,\kappa\cdot\max_{T\cap\mathrm{slab}}
  \frac{h_T}{\sqrt{w_T^Q}}
  +\sqrt{\frac{H_t|R|Y}{3}},
\]
with $|R|=\tfrac12$.  Any $H_t$ is admissible; the code optimizes $H_t$
numerically (correctness-free).  The certified level-one run used this
bound and therefore stands a fortiori under $\tau_h=0$.

\section{Certificate tables}

\begin{table}[ht]
\centering
\begin{tabular}{@{}lcc@{}}
\toprule
Quantity & Arb upper bound & Role \\
\midrule
$\gamma$ & $0.09847$ & mass-error / stiffness \\
$\tau$ & $0.29085$ & historical slab trace error \\
$\alpha_h$ & $0.02240$ & volume-functional error \\
$S_M$ & $1.34\cdot 10^{-1}$ & sliver mass \\
$S_Q$ & $7.45\cdot 10^{-3}$ & sliver stiffness \\
$S_\Sigma$ & $5.71\cdot 10^{-3}$ & sliver floor stiffness \\
$V_\Sigma$ & $6.94\cdot 10^{-3}$ & sliver weighted volume \\
$\|\ell_{\mathrm{rad}}\|$ & $6.6\cdot 10^{-8}$ & $\ell_0$ enclosure radius \\
$\|Q_r\|$ & $2.5\cdot 10^{-13}$ & stiffness radius \\
Rump shift $c$ & $\approx 2\cdot 10^{-8}$ & PSD certificate \\
Float margin & $\approx 10^{-4}$ & diagnostic only \\
\bottomrule
\end{tabular}
\caption{Abridged certificate constants for Theorem~\ref{thm:main}
(2026-07-10).  All eight $s$-windows: $c_e>d_e$ and Rump PSD true.}
\label{tab:cert1}
\end{table}

\begin{thebibliography}{99}

\bibitem{BB11}
V.~Blomer and F.~Brumley,
\emph{On the Ramanujan conjecture over number fields},
Ann.\ of Math.\ (2) \textbf{174} (2011), 581--605.

\bibitem{BSV06}
A.~R.~Booker, A.~Str\"ombergsson, and A.~Venkatesh,
\emph{Effective computation of Maass cusp forms},
Int.\ Math.\ Res.\ Not.\ 2006, Art.\ ID 71281.

\bibitem{Bebendorf03}
M.~Bebendorf,
\emph{A note on the Poincar\'e inequality for convex domains},
Z.\ Anal.\ Anwend.\ \textbf{22} (2003), 751--756.

\bibitem{CG14}
C.~Carstensen and J.~Gedicke,
\emph{Guaranteed lower bounds for eigenvalues},
Math.\ Comp.\ \textbf{83} (2014), 2605--2629.

\bibitem{CDMN26}
C.~Carstensen, A.~Dond, N.~Maity, and N.~Nataraj,
\emph{Lower energy bounds for the Stokes eigenvalue problem},
Comput.\ Math.\ Appl.\ \textbf{214} (2026), 27--50.

\bibitem{CP22}
C.~Carstensen and S.~Puttkammer,
\emph{Adaptive guaranteed lower eigenvalue bounds with optimal
convergence rates},
arXiv:2203.01028.

\bibitem{CZZ20}
C.~Carstensen, Z.~Zhai, and S.~Zhang,
\emph{A framework for the error analysis of discontinuous finite element
methods for elliptic eigenvalue problems and applications},
SIAM J.\ Numer.\ Anal.\ \textbf{58} (2020), 109--124.

\bibitem{EGM98}
J.~Elstrodt, F.~Grunewald, and J.~Mennicke,
\emph{Groups Acting on Hyperbolic Space},
Springer, 1998.

\bibitem{Friedman}
J.~Friedman,
\emph{The Selberg trace formula for $\mathrm{PSL}(2,O_K)$},
arXiv:math/0612807.

\bibitem{Leutbecher67}
A.~Leutbecher,
\emph{\"Uber die Heckeschen Gruppen $G(\lambda)$},
Math.\ Z.\ \textbf{100} (1967), 183--200.

\bibitem{PaperII}
\emph{Exceptional eigenvalues of Bianchi orbifolds by verified finite
elements.  II: Congruence levels of the Picard group},
companion paper (this series).

\bibitem{PaperIII}
\emph{Exceptional eigenvalues of Bianchi orbifolds by verified finite
elements.  III: The Eisenstein--Picard orbifold and its first
congruence cover},
companion paper (this series).

\bibitem{PW60}
L.~E.~Payne and H.~F.~Weinberger,
\emph{An optimal Poincar\'e inequality for convex domains},
Arch.\ Rational Mech.\ Anal.\ \textbf{5} (1960), 286--292.

\bibitem{Rump06}
S.~M.~Rump,
\emph{Verification of positive definiteness},
BIT Numer.\ Math.\ \textbf{46} (2006), 433--452.

\bibitem{Shimizu63}
H.~Shimizu,
\emph{On discontinuous groups operating on the product of the upper
half planes},
Ann.\ of Math.\ (2) \textbf{77} (1963), 33--71.

\end{thebibliography}

\end{document}
